\documentclass{tietreport}

\setlength{\headheight}{14.5pt} % fixes fancyhdr warning with 12pt font

\emergencystretch=2em

% Import images
\usepackage{graphicx}

% Extra packages used for tables and rotated (landscape) diagrams
\usepackage{booktabs}
\usepackage{tabularx}
\usepackage{array}
\usepackage{rotating}
\usepackage{float}
\usepackage{enumitem}
\setlist{itemsep=2pt,topsep=4pt,parsep=0pt}
\makeatletter\@openrightfalse\makeatother % no blank pages before chapters
\renewcommand{\topfraction}{0.95}
\renewcommand{\bottomfraction}{0.95}
\renewcommand{\textfraction}{0.05}
\renewcommand{\floatpagefraction}{0.7}
\setlength{\textfloatsep}{10pt plus 2pt minus 2pt}
\setlength{\intextsep}{10pt plus 2pt minus 2pt}

% Bibliography configuration
\usepackage[
  date=year,
  style=alphabetic,
  backend=biber,
  sorting=nty,
  maxnames=5,
  minnames=3,
  maxalphanames=4,
  minalphanames=3,
  backref=true,
  doi=false,
  isbn=false,
  url=false,
  eprint=false
]{biblatex}
\addbibresource{references.bib}

%% Helpers ------------------------------------------------------------
% Wide diagram on its own rotated page: file, caption, label
\newcommand{\widefig}[3]{%
  \begin{figure}[!htbp]
    \centering
    \makebox[\textwidth][c]{\includegraphics[width=1.12\textwidth]{figures/#1}}
    \caption{#2}\label{#3}
  \end{figure}}
% Square diagram in the normal text flow: file, caption, label
\newcommand{\squarefig}[3]{%
  \begin{figure}[!htbp]
    \centering
    \includegraphics[width=0.58\textwidth]{figures/#1}
    \caption{#2}\label{#3}
  \end{figure}}
\newcolumntype{Y}{>{\raggedright\arraybackslash}X}

%% ==================================================
%% CUSTOMIZE THESE FIELDS FOR YOUR PROJECT
%% ==================================================

% Course/Lab header
\TitlePageHeader{%
  UCS503P - Software Engineering%
}

% Main project title
\ProjectTitle{%
  ArogyaMitra%
}

% Subtitle or project description
\ProjectSubTitle{%
  An AI-Assisted Real-Time Telemedicine Consultation Platform%
}

% Student information (up to 4 students)
\author{%
  \texttt{1024170375} Animesh Sudhanshu \\
  \texttt{1024170372} Ashish Bhagat \\
  \texttt{1024170380} Devansh Rathaur%
}

% Group number, degree, year, program
\TitlePageSubText{%
  Group: \texttt{3Q34} BE Third Year, CSE%
}

% Faculty advisor/guide name
\AdvisorName{%
  Dr. Paramveer Kaur%
}

% Evaluation period and department info
\TitlePageFooterText{{%
    \large Mid-Semester Evaluation \\[0.2cm]
    \bfseries Computer Science and Engineering
    Department%
  } \\[0.15cm]
  Thapar Institute of Engineering and Technology,
  Patiala%
}

%% ==================================================

\begin{document}

% Generate title page
\maketitle

% Table of contents (auto-generated from chapters)
\setcounter{tocdepth}{1}
\tableofcontents

% ===== CHAPTER 1: INTRODUCTION =====

\chapter{Introduction}

\section{Background and Context}

ArogyaMitra is a web-based, AI-assisted real-time telemedicine
platform. It reduces the friction between a patient who needs
medical advice and a doctor who can give it, for non-emergency
consultations, using chat, online booking and an AI
assistant that prepares context before the doctor joins.

The broader motivation for the project is improving access to
primary healthcare and healthcare equity. The immediate engineering
objective is to translate low-friction, real-time doctor--patient
consultation into a measurable, deployable software solution built
and iterated on by a small student team. Real-time web
technologies and affordable managed backends have matured enough that a
small team can build a credible consultation platform in an academic
timeframe, while conversational AI can plausibly assist (not
replace) the pre-consultation triage step.

This is an engineering course project, so the focus is on
integration, workflow correctness and safe deployment of existing,
well-understood technologies rather than on novel algorithm
research. Time-to-value is the primary heuristic: meaningful
increments are delivered quickly by adapting an existing open-source
reference implementation, prototyping the core consultation workflow
and validating it within weekly iterations, with CI-backed
automation introduced early to reduce release friction.

\subsection{Problem Statement}

Patients, particularly those in smaller towns with limited
specialist availability, face delays and information loss when
seeking timely medical advice for non-emergency conditions. The
existing challenges are:

\begin{itemize}
  \item \textbf{Access friction:} Booking an in-person visit means
    travel and waiting, even for simple, non-emergency queries.
  \item \textbf{Fragmented communication:} Follow-ups, advice and
    prescriptions are exchanged informally (calls, chat apps) with
    no persistent record.
  \item \textbf{No triage before consultation:} Patients cannot
    get a first-pass sense of urgency or prepare context before a
    doctor is available.
  \item \textbf{Scheduling opacity:} Patients cannot easily see
    doctor availability or track the status of a booked
    consultation.
\end{itemize}

\section{Project Objectives}

ArogyaMitra streamlines the consultation journey end to end through
the following objectives:

\begin{enumerate}
  \item \textbf{Patient portal and doctor dashboard:}
    Registration, profiles, doctor discovery, booking, consultation
    queue and history.
  \item \textbf{Real-time chat consultation:} A secure chat room
    per consultation, joinable only in the booked time window.
  \item \textbf{AI pre-consultation assistant:} Gemini-based symptom
    intake with an emergency-triage safety layer; the assistant
    never diagnoses or prescribes.
  \item \textbf{Records and admin control:} Medical records and
    admin approval of doctors.
\end{enumerate}

\noindent\textbf{Current status.} The frontend and backend currently
run locally on the development machine (localhost) and have not yet
been deployed to Netlify and Render. Online payments (Stripe) and
video conferencing are not part of the current implementation and
will be implemented in the future.

The \textbf{primary success metric} is a median time-to-connect
(TTC) of at most 2 minutes between the scheduled start and both
parties being in the same session.

\section{Feasibility Study}

Table~\ref{tab:feasibility} answers the six feasibility questions
for the proposed platform.

\begin{table}[!htbp]
\centering
\small
\renewcommand{\arraystretch}{1.25}
\begin{tabularx}{\textwidth}{@{}l>{\raggedright\arraybackslash}p{1.35in}Y@{}}
  \toprule
  \textbf{No.} & \textbf{Question} & \textbf{Answer} \\
  \midrule
  1 & What if the system wasn't implemented? &
    Patients keep travelling and waiting for simple queries.
    Advice stays scattered over calls and chat apps with no
    record. \\
  2 & What are the current process problems? &
    No instant view of doctor availability or booking status.
    No triage or context before the consultation starts. \\
  3 & How will the proposed system help? &
    Remote booking, chat, persistent history and records.
    AI prepares symptoms for the doctor; emergencies are
    escalated. \\
  4 & What will be the integration problems? &
    Join control by booked hour. Clerk-JWT secured API; keeping AI
    vectors in sync via hooks. \\
    hooks. \\
  5 & Is new technology needed? Which skills? &
    No new hardware: Next.js, Node/Express, MongoDB. Skills:
    full-stack JavaScript, REST, prompt/tool design. \\
  6 & What facilities must be supported? &
    Managed hosting (Netlify, Render, planned; currently
    localhost) and MongoDB Atlas. SaaS: Clerk and Gemini; a
    browser. \\
  \bottomrule
\end{tabularx}
\caption{Feasibility study of the ArogyaMitra platform}
\label{tab:feasibility}
\end{table}

\noindent\textbf{Feasibility verdict:} Highly viable. The platform
uses a standard managed stack, is built on a proven open-source
codebase, and has minimal infrastructure demands.

\section{Scope and Deliverables}

\subsection{Initial Deliverable (First Iteration)}

\begin{itemize}
  \item Fork/clone the reference repository and set up the local
    environment (done; the staging environment is planned).
  \item Patient and doctor authentication and profile pages.
  \item Appointment booking flow with doctor availability.
  \item Real-time chat consultation between doctor and patient.
  \item Basic logging and timestamps for TTC measurement.
  \item CI: build, lint and backend unit tests.
  \item CD to staging (Render and Netlify) with a single pipeline
    job (planned; the application currently runs on localhost).
\end{itemize}

\subsection{Subsequent Deliverables}

\begin{itemize}
  \item Online payments (Stripe) and real-time video consultation
    (LiveKit): to be implemented in the future.
  \item Deployment of the frontend (Netlify) and backend (Render)
    from the current local setup.
  \item AI pre-consultation assistant with the emergency-triage
    safety layer.
  \item In-app alerts for upcoming consultations.
  \item Admin/analytics view of bookings, TTC and completion rate.
  \item Regression test suite for AI assistant safety behaviour.
\end{itemize}

% ===== CHAPTER 2: METHODOLOGY & DESIGN =====

\chapter{Methodology and Design}

\section{Requirements Engineering}

\subsection{Requirements Engineering Process}

The requirements were developed through the seven standard
requirements-engineering activities: (1)~inception,
(2)~elicitation, (3)~elaboration, (4)~negotiation,
(5)~specification, (6)~validation and (7)~management.

\subsection{Stakeholders}

\begin{itemize}
  \item \textbf{Primary actors:} Patient and Doctor (end users).
  \item \textbf{Secondary actor:} Administrator (approval and
    account management).
  \item \textbf{External systems:} Clerk and Gemini; Stripe and
    LiveKit will be added in the future.
\end{itemize}

\subsection{Functional Requirements}

\begin{table}[!htbp]
\centering
\small
\renewcommand{\arraystretch}{1.25}
\begin{tabularx}{\textwidth}{@{}l>{\raggedright\arraybackslash}p{1.45in}Y@{}}
  \toprule
  \textbf{ID} & \textbf{Requirement} & \textbf{Description} \\
  \midrule
  FR1 & User authentication and roles &
    Clerk sign-in with JWT sessions; patient, doctor and admin
    roles. \\
  FR2 & Doctor onboarding and approval &
    Doctors publish specialty, price and schedule; the admin
    approves or rejects. \\
  FR3 & Discovery and booking &
    Patients search doctors and book a date and time slot. \\
  FR4 & Payments &
    Online payments for paid doctors (Stripe Checkout) will be
    implemented in the future. \\
  FR5 & Real-time consultation &
    Text chat in a per-consultation room, joinable only in the
    booked hour; video calls to be added in the future. \\
  FR6 & Consultation lifecycle &
    pending $\rightarrow$ in-progress $\rightarrow$ completed /
    canceled; hourly auto-close. \\
  FR7 & AI pre-consultation assistant &
    Symptom triage tools, emergency override, RAG over the
    patient's own records, PDF upload, hourly quota. \\
  FR8 & Records and alerts &
    Consultation history, medical records and in-app consultation
    alerts. \\
  \bottomrule
\end{tabularx}
\caption{Functional requirements}
\label{tab:fr}
\end{table}

\subsection{Non-Functional Requirements}

\begin{table}[!htbp]
\centering
\small
\renewcommand{\arraystretch}{1.25}
\begin{tabularx}{\textwidth}{@{}l>{\raggedright\arraybackslash}p{1.2in}Y@{}}
  \toprule
  \textbf{ID} & \textbf{Category} & \textbf{Requirement} \\
  \midrule
  NFR1 & Performance &
    Median time-to-connect $\leq$ 2 min; streamed AI replies. \\
  NFR2 & Security &
    Clerk JWT, role checks, participant check, helmet, input
    sanitising, rate limiting. \\
  NFR3 & Usability &
    Responsive web UI for desktop, tablet and low-to-mid range
    phones. \\
  NFR4 & Reliability &
    AI context search degrades gracefully. \\
  NFR5 & Privacy and safety &
    Per-patient vector scope; the AI never diagnoses or
    prescribes. \\
  \bottomrule
\end{tabularx}
\caption{Non-functional requirements}
\label{tab:nfr}
\end{table}

\section{System Architecture}

\subsection{High-Level Design}

ArogyaMitra follows a decoupled three-tier web architecture with
external SaaS services (Figure~\ref{fig:architecture}). The client
tier is a responsive web interface (desktop, tablet and phone
browsers) with patient, doctor and admin portals. The presentation
tier is a Next.js~16 / React~18 application, and the application
tier is a Node.js/Express REST API (\texttt{/v1}). The target
deployment is Netlify for the frontend and Render for the backend,
but both currently run locally on the development machine
(localhost). The data tier is MongoDB Atlas accessed through the
Mongoose ODM. The tiers communicate over HTTPS, REST with a Clerk
JWT, and Mongoose respectively. Clerk and the Google Gemini API are
consumed as external services. Payment (Stripe) and video (LiveKit)
services shown in the architecture figure are not part of the
current implementation and will be integrated in the future.

\widefig{architecture.png}{System architecture: three-tier web
  architecture with external SaaS services}{fig:architecture}

\subsection{Technical Stack}

\begin{itemize}
  \item \textbf{Framework/Language:} Next.js and React with
    Tailwind and shadcn/ui on the frontend~\cite{nextjs}; Node.js
    and Express on the backend~\cite{expressjs}.
  \item \textbf{Database:} MongoDB Atlas with the Mongoose ODM and
    Atlas Vector Search for patient-scoped embeddings~\cite{mongodbatlas}.
  \item \textbf{Tools:} GitHub Actions for CI (build, lint and
    tests); Render (backend) and Netlify (frontend) for staging
    deployment, which is planned (the application currently runs on
    localhost).
  \item \textbf{Libraries and services:} Clerk for authentication
    and JWT sessions~\cite{clerk}; Google Gemini (via the Vercel AI
    SDK) for the AI assistant~\cite{gemini}.
  \item \textbf{Planned for the future:} Stripe for online payments
    and LiveKit for video consultation.
\end{itemize}

\subsection{AI Assistant Scope}

The AI pre-consultation assistant is an integration task, not a
research contribution:

\begin{itemize}
  \item An existing LLM API (Gemini) is used for conversational
    symptom intake; there is no in-house model training.
  \item A lightweight safety/triage layer detects emergency-flagged
    symptoms (for example chest pain, stroke signs, self-harm) and
    always routes to a ``seek immediate care'' response rather than
    open-ended advice.
  \item The assistant only prepares context for the doctor; it
    never issues a diagnosis or prescription.
\end{itemize}

\subsection{Scalability and Deployability}

\begin{itemize}
  \item \textbf{Theoretical scalability:} The frontend and backend
    are decoupled; REST endpoints are stateless; appointment and
    user queries are indexed in MongoDB.
  \item \textbf{Practical deployability:} The system is designed to
    run on common managed hosting: MongoDB Atlas as the database,
    Render and Netlify as deployment targets, and GitHub Actions for
    the staging/production pipeline. Deployment is still to be done;
    at present everything runs locally.
\end{itemize}

\section{UML Modelling}

The diagrams in this section describe the target design of the
platform. Elements for online payments (Stripe) and video
conferencing (LiveKit) that appear in them will be implemented in
the future.

\subsection{Use Case Model}

The use case model (Figure~\ref{fig:usecase}) identifies the
Patient, Doctor and Admin as human actors, and Clerk, Gemini and a
cron scheduler as external systems in the current scope. The Stripe
and LiveKit actors shown in the diagram stand for online payments
and video conferencing, which will be implemented in the future.

\widefig{usecase.png}{Use case diagram of the ArogyaMitra
  platform}{fig:usecase}

\subsection{Class Diagrams}

The class diagrams are derived from the Mongoose models of the
backend. Figure~\ref{fig:class-domain} shows the domain model (User
and its Patient, Doctor and Admin specialisations, Consultation and
Payment, and their enumerations). Figure~\ref{fig:class-ai} shows the
AI assistant and medical-records classes (AiConversation, AiUsage,
MedicalEmbedding and the RAG path).

\widefig{class_domain.png}{Class diagram: domain model}{fig:class-domain}

\widefig{class_ai.png}{Class diagram: AI assistant and medical records
  (AiConversation, AiUsage, MedicalEmbedding, RAG
  path)}{fig:class-ai}

\subsection{Activity Diagram}

Figure~\ref{fig:activity} shows the end-to-end workflow, from
sign-in to a completed consultation.

\widefig{activity.png}{Activity diagram: end-to-end workflow from
  sign-in to completed consultation}{fig:activity}

\subsection{Swimlane Diagram}

Figure~\ref{fig:swimlane} shows who does what across the Patient,
Platform, External services and Doctor lanes.

\widefig{swimlane.png}{Swimlane diagram: Patient, Platform, External
  services and Doctor}{fig:swimlane}

\subsection{State Diagrams}

Figure~\ref{fig:states} shows the Consultation, Payment and
Doctor-approval lifecycles. A consultation is created as
\emph{pending}, becomes \emph{in-progress} when a participant joins
within the booked hour, and ends as \emph{completed} or
\emph{canceled}; a cron job runs at minute~0 of every hour to close
past consultations.

\widefig{states.png}{State diagrams: Consultation, Payment and
  Doctor-approval lifecycles}{fig:states}

\subsection{Sequence Diagrams}

One sequence diagram describes the key AI interaction: the
pre-consultation flow of quota check, Gemini tool calls,
triage/RAG and persistence (Figure~\ref{fig:seq-ai}).

\widefig{seq_ai.png}{Sequence diagram: AI pre-consultation assistant
  (quota, Gemini tool calls, triage/RAG, persistence)}{fig:seq-ai}

\section{Data Flow Diagrams}

The data flow diagrams use the following notation: a circle is a
process, a rectangle is an external entity, and an open box is a
data store. Figure~\ref{fig:dfd0} gives the Level~0 context
diagram, Figure~\ref{fig:dfd1} the Level~1 decomposition, and
Figures~\ref{fig:dfd2-1}--\ref{fig:dfd2-6} the Level~2 diagrams of
processes 1.0 to 6.0.

\widefig{dfd0.png}{DFD Level 0: context diagram}{fig:dfd0}

\widefig{dfd1.png}{DFD Level 1}{fig:dfd1}

\squarefig{dfd2_1.png}{DFD Level 2: process 1.0, Authentication and
  user sync}{fig:dfd2-1}

\squarefig{dfd2_2.png}{DFD Level 2: process 2.0, Doctor profile and
  approval}{fig:dfd2-2}

\squarefig{dfd2_3.png}{DFD Level 2: process 3.0, Booking and
  payment}{fig:dfd2-3}

\squarefig{dfd2_4.png}{DFD Level 2: process 4.0, Consultation
  session}{fig:dfd2-4}

\squarefig{dfd2_5.png}{DFD Level 2: process 5.0, AI pre-consultation
  assistant}{fig:dfd2-5}

\squarefig{dfd2_6.png}{DFD Level 2: process 6.0, Medical records and
  embeddings}{fig:dfd2-6}

\section{SDLC Model Selection}

Waterfall, Prototyping and Agile Scrum were evaluated
(Table~\ref{tab:sdlc}); Agile Scrum was selected.

\begin{table}[!htbp]
\centering
\footnotesize
\renewcommand{\arraystretch}{1.25}
\begin{tabularx}{\textwidth}{@{}>{\raggedright\arraybackslash}p{0.85in}YYY@{}}
  \toprule
  \textbf{Metric} & \textbf{Waterfall} & \textbf{Prototyping} &
  \textbf{Agile Scrum (selected)} \\
  \midrule
  Requirement adaptability &
    Rigid; changes are expensive once a phase concludes &
    Flexible for UI; weak for architectural evolution &
    High; requirements refined iteratively every sprint \\
  Feedback frequency &
    Late; user evaluation only at final acceptance &
    Early UI feedback; variable back-end validation &
    Continuous; pilot users and supervisor review after each
    sprint \\
  Risk and bug mitigation &
    Found late, during integrated testing &
    Medium; addressed during prototype refinement &
    Immediate; CI (build + lint + tests) on every change, AI
    safety regression suite \\
  Delivery milestones &
    Single monolithic release at the end date &
    Demo prototype followed by a full build &
    Incremental; usable MVP first (auth + booking + chat), then
    AI and admin \\
  \bottomrule
\end{tabularx}
\caption{Evaluation of SDLC models}
\label{tab:sdlc}
\end{table}

\begin{itemize}
  \item \textbf{How Scrum is applied:} Weekly iterations, each
    ending with a working increment. CI runs on GitHub Actions;
    staging on Render and Netlify is planned, since the
    application currently runs on localhost.
  \item \textbf{Release plan:} Sprint~1 delivers auth, profiles,
    booking, chat and TTC logging. Next come the AI assistant with
    safety tests, admin and deployment. Online payments and video
    consultation are planned for the future.
  \item \textbf{Why not the others:} Waterfall: AI and real-time
    behaviour need early feedback. Prototyping: weak for back-end
    integrity.
\end{itemize}

Agile Scrum supports iterative refinement, continuous feedback,
early delivery and rapid risk mitigation.

% ===== CHAPTER 3: RESULTS & EVALUATION =====

\chapter{Results and Evaluation}

\section{Testing Methodology}

The frontend and backend currently run on the local development
machine (localhost) and have not yet been deployed to Netlify and
Render. This chapter first records the current implementation
status and then defines the evaluation plan against which the
implementation will be measured.

\subsection{Current Implementation Status}

\begin{itemize}
  \item \textbf{Environment:} Frontend and backend run locally on
    the development PC (localhost). Deployment to Netlify and Render
    is planned.
  \item \textbf{Payments and video:} The Stripe payment gateway and
    video conferencing are not implemented yet and will be
    implemented in the future.
\end{itemize}

\subsection{Testing Plan}

\begin{itemize}
  \item \textbf{Unit Testing:} Backend unit tests run in CI
    together with the build and lint steps on every change.
  \item \textbf{Integration Testing:} Currently performed on
    localhost. Once deployed, repeatable deployments to the staging
    environment (Render and Netlify) via the CI/CD pipeline will
    exercise login, booking and consultation end to end.
  \item \textbf{Performance Testing:} Time-to-connect is measured
    from server-side timestamps (scheduled time versus
    session-join event).
  \item \textbf{User Testing:} A pilot with 10--15 patients and
    3--5 volunteer or simulated doctor accounts, comparing TTC and
    booking completion rate across iterations.
  \item \textbf{AI Safety Testing:} A fixed regression suite of
    symptom prompts, including emergency cases, is run before each
    release to catch safety regressions.
\end{itemize}

\section{Performance Metrics}

Table~\ref{tab:performance} lists the metrics and targets used to evaluate ArogyaMitra.

\begin{table}[H]
\centering
\small
\renewcommand{\arraystretch}{1.25}
\begin{tabularx}{\textwidth}{@{}>{\raggedright\arraybackslash}p{1.3in}>{\raggedright\arraybackslash}p{1.45in}Y@{}}
  \toprule
  \textbf{Metric} & \textbf{Target} & \textbf{How it is measured} \\
  \midrule
  Time-to-connect (primary) & Median $\leq$ 2 minutes &
    Server-side timestamps: scheduled start versus both parties
    in the same session \\
  Booking completion rate & Compared across iterations &
    Share of started bookings that result in a confirmed
    appointment \\
  AI assistant safety accuracy & Fixed suite must pass before every release &
    Share of emergency-flagged test cases routed to ``seek
    immediate care'' \\
  User clarity & Simple 1--5 rating &
    Patient-reported clarity of the consultation flow \\
  Reliability & e.g., $\geq$ 99\% availability &
    Availability of login, booking and consultation on
    staging/production during test windows \\
  \bottomrule
\end{tabularx}
\caption{Evaluation metrics and targets (to be measured during the pilot)}
\label{tab:performance}
\end{table}

\section{Analysis}

\begin{itemize}
  \item \textbf{Objectives:} Feasibility has been validated, the
    requirements are documented, and the UML and DFD models are
    complete. Online payments (Stripe), video conferencing and the
    Netlify/Render deployment are not yet implemented and will be
    done in the future; measured results will follow once they are
    done.
  \item \textbf{Challenges and mitigations:}
  \begin{itemize}
    \item \emph{Data privacy and consent:} collect only the
      medical data required for the consultation record, and do
      not persist raw AI conversations beyond what is needed for
      the doctor's context.
    \item \emph{AI safety failure:} a fixed emergency-prompt
      regression suite must pass before every release; the
      assistant never diagnoses or prescribes.
    \item \emph{Real-time reliability:} connection-state handling
      and reconnect logic so that a dropped session does not lose
      consultation context.
    \item \emph{Doctor adoption:} keep the doctor dashboard minimal
      to reduce onboarding friction during the pilot.
    \item \emph{Operational issues in the pilot:} deploy to staging
      early (the application currently runs only on localhost) and
      ensure CI/CD produces repeatable deployments before the
      pilot window.

  \end{itemize}
\end{itemize}

% ===== CHAPTER 4: CONCLUSION =====

\chapter{Conclusion}

\section{Summary of Work}

This project proposes a deployable, AI-assisted real-time
telemedicine platform that reduces the friction between patients
and doctors for non-emergency consultations. It emphasises
time-to-value through rapid iterations on a proven codebase, CI/CD
for frequent safe releases, and measurable evaluation criteria
(time-to-connect and AI assistant safety accuracy), staying focused
on engineering integration and safe AI-assistance boundaries rather
than research-driven novelty. At present the frontend and backend
run locally on the development machine. Online payments (Stripe),
video conferencing and the deployment to Netlify and Render are not
part of the current implementation and will be implemented in the
future.

\section{Key Findings}

\begin{itemize}
  \item \textbf{Feasibility validated:} The platform is highly
    viable on a standard managed stack with minimal infrastructure
    demands.
  \item \textbf{SRS documented:} Eight functional and five
    non-functional requirements were specified.
  \item \textbf{UML and DFD modelled:} UML diagrams and DFD
    Levels~0--2 were produced.
  \item \textbf{Agile Scrum selected:} It best supports iterative
    refinement, continuous feedback, early delivery and rapid risk
    mitigation.
\end{itemize}

\section{Future Work and Recommendations}

\begin{enumerate}
  \item Implement online payments with the Stripe payment gateway
    (checkout, webhook and idempotent fulfilment) in a future
    iteration.
  \item Implement real-time video consultation (LiveKit) alongside
    the existing chat in a future iteration.
  \item Deploy the frontend to Netlify and the backend to Render,
    with CI/CD to staging, replacing the current localhost setup.
  \item Add the AI pre-consultation assistant with its
    emergency-triage regression suite and the admin/analytics view,
    then run the pilot and compare TTC and booking completion rate
    across iterations.
\end{enumerate}

\section{Final Remarks}

The specification, models and process selected in this report give
the team a clear and validated basis for building ArogyaMitra
incrementally, with safety boundaries for the AI assistant defined
from the start.

% ===== BIBLIOGRAPHY =====

% Print bibliography with heading in TOC
\printbibliography[heading=bibintoc]

\end{document}
